% Created 2022-03-14 ma 16:23
% Intended LaTeX compiler: pdflatex
\documentclass[a4paper]{article}
\usepackage[utf8]{inputenc}
\usepackage[T1]{fontenc}
\usepackage{graphicx}
\usepackage{longtable}
\usepackage{wrapfig}
\usepackage{rotating}
\usepackage[normalem]{ulem}
\usepackage{amsmath}
\usepackage{amssymb}
\usepackage{capt-of}
\usepackage{hyperref}
\usepackage{minted}
\usepackage{a4wide}
\usepackage{minted}
\setminted[python]{linenos=true}
\setminted[python]{frame=lines}
\usepackage{fouriernc}
\renewcommand{\P}[1]{\,\mathsf{P}\left[#1\right]}
\newcommand{\E}[1]{\,\mathsf{E}\left[#1\right]}
\author{Nicky van Foreest}
\date{2022-03-14}
\title{Memoization with python}
\begin{document}

\maketitle
\tableofcontents

\(\def\RR{{\bf R}}
    \def\bold#1{{\bf #1}}\)


\section{Content of a Discrete Hyper Pyramid}
\label{sec:orgb3281fa}

We like to compute the number of possibilities \(\P{n, N}\) for \(x = (x_1, \ldots x_n)\) such that \(x_i \in \{0,1,\ldots, N\}\) and \(\sum_i x_i \leq N\).
It is easy to see that \(\P{n, N}\) satisfies the recursion:
$$\P{n, N} = \sum_{i=0}^N \P{n-1, N-i},$$
with boundary conditions \(\P{1, N} = N+1\) for all \(N\).
Note that by using the summation above, this condition can be replaced by \(\P{0, N} = 1\) for all \(N\).

Computing \(\P{n,N}\) is easy when we use  memoization. In fact,  we can code the compuation in nearly one line!

\begin{minted}[]{python}
from functools import lru_cache

@lru_cache(maxsize=128)
def P(n, N):
   return (n==0) or sum( P(n-1, N-i) for i in range(N+1) )

n=5
N=80
print(P(n = 5, N = 80))
\end{minted}

\begin{verbatim}
32801517
\end{verbatim}

\section{A probability problem}
\label{sec:orga114cf3}

We throw multiple times with a coin that lands heads with probability \(p\).
What is the probability \(\P{n,k}\) to see at least \(k\) heads in row when you throw a coin \(n\) times?

A bit of thinking shows  that \(\P{n,k}\) must satisfy the recursion
$$\P{n,k} = p^k + \sum_{i=1}^k p^{i-1} q\, \P{n-i,k},$$
because it is possible to throw \(k\) times heads from the first throw, but otherwise you throw \(i\), \(i<k\), times a heads, then a tails, after which you have to start all over again.

Reasoning similarly, the expected number times \(\E{n,k}\) to see at least \(k\) heads in row when you throw a coin \(n\) times must satisfy the recursion
$$\E{n,k} = p^k(1+\E{n-k,k}) + \sum_{i=1}^k p^{i-1} q\, \E{n-i,k}.$$


\begin{minted}[]{python}
p = 0.5
q = 1. - p

@lru_cache(maxsize=128)
def P(n, k):
    """
    probability to see at least k heads in row when a coin is thrown n times
    """
    if n < k:
        return 0
    else:
        return sum(P(n-i,k) * p**(i-1) * q for i in range(1,k+1)) + p**k

@lru_cache(maxsize=128)
def E(n, k):
    """
    expected number of times to see at least k heads in row when a coin is thrown n times
    """
    if n < k:
        return 0
    else:
        tot = sum(E(n-i,k) * p**(i-1) *q for i in range(1,k+1))
        tot += p**k * (1 + E(n-k,k))
        return tot



k = 2

for n in range(k,10):
    print(n, P(n,k), E(n,k))
\end{minted}

\begin{verbatim}
2 0.25 0.25
3 0.375 0.375
4 0.5 0.5625
5 0.59375 0.71875
6 0.671875 0.890625
7 0.734375 1.0546875
8 0.78515625 1.22265625
9 0.826171875 1.388671875
\end{verbatim}
\end{document}